%%%%%%%%%%  scan idx 215  |  folio 202  %%%%%%%%%%
dans $F$ d'ordre $\le n$.

\textbf{1.2.3.} Nous aurons aussi à considérer des opérateurs différentiels à coefficients complexes. Notons $\cO_S\otimes\mathbb C$ (resp. $\cO_X\otimes\mathbb C$) le complexifié de $\cO_S$ (resp. $\cO_X$). Soient $E,F$ des $(\cO_X\otimes\mathbb C)$-Modules localement libres de type fini. Par un opérateur différentiel de $E$ dans $F$ d'ordre $\le n$, on entendra un homomorphisme $f^{-1}(\cO_S\otimes\mathbb C)$-linéaire de $E$ dans $F$ qui se factorise à travers $d^n_{X/S}(E)$ en un homomorphisme $(\cO_X\otimes\mathbb C)$-linéaire de $P^n_{X/S}(E)$ dans $F$. Si l'on désigne par $\Diff^n_\mathbb C(E,F)$ l'ensemble de ces opérateurs, l'opérateur $d^n_{X/S}(E)$ définit donc une bijection
\begin{equation}\tag{1.2.3.1}
  \Homop_{\cO_X\otimes\mathbb C}\bigl(P^n_{X/S}(E),F\bigr)
  \xrightarrow{\sim}
  \Diff^n_\mathbb C(E,F).
\end{equation}

\textbf{1.2.4.} Pour $n\ge1$, on a une suite exacte canonique
\[
  0\longrightarrow S^n(\Omega^1_{X/S})
  \longrightarrow P^n_{X/S}
  \longrightarrow P^{n-1}_{X/S}
  \longrightarrow 0.\text{\footnotemark}
\]
\footnotetext{$S^n(\ )$ désigne la puissance symétrique $n$-ième.}%
D'où, en tensorisant sur $\cO_X$\footnote{Pour la deuxième structure de $\cO_X$-Module de $P^n_{X/S}$ et $P^{n-1}_{X/S}$.} par un $\cO_X$-Module (resp. $(\cO_X\otimes\mathbb C)$-Module) localement libre de type fini $E$, une suite exacte de $\cO_X$ (resp. $(\cO_X\otimes\mathbb C)$)-Modules~:
\begin{equation}\tag{1.2.4.1}
  0\longrightarrow S^n(\Omega^1_{X/S})\otimes_{\cO_X}E
  \longrightarrow P^n_{X/S}(E)
  \longrightarrow P^{n-1}_{X/S}(E)
  \longrightarrow 0.
\end{equation}
Soit $F$ un $\cO_X$ (resp. $(\cO_X\otimes\mathbb C)$)-Module localement libre de type fini. Appliquant à (1.2.4.1) le foncteur $\Homop_{\cO_X}(.,F)$ (resp. $\Homop_{\cO_X\otimes\mathbb C}(.,F)$), on en déduit, compte tenu des identifications (1.2.2.1) et (1.2.3.1), une suite exacte
\begin{equation}\tag{1.2.4.2}
\begin{gathered}
  0\longrightarrow \Diff^{n-1}(E,F)
  \longrightarrow \Diff^n(E,F)
  \xrightarrow{\sigma_n}
  \Homop_{\cO_X}\bigl(S^n(\Omega^1_{X/S})\otimes_{\cO_X}E,F\bigr)\\
  (\text{resp. } 0\longrightarrow \Diff^{n-1}_\mathbb C(E,F)
  \longrightarrow \Diff^n_\mathbb C(E,F)
  \longrightarrow
  \Homop_{\cO_X\otimes\mathbb C}\bigl(S^n(\Omega^1_{X/S})\otimes_{\cO_X}E,F\bigr)).
\end{gathered}
\end{equation}
